

\section{Introduction}   \label{ss:intro}

Practical computation for statistical mechanical models is often built around transfer matrix methods (see e.g. \cite{Baxter81,Baxter83,LiebMattis,Martin91,Wannier,Brown67} %,other} 
and references therein). 
Transfer matrix methods in turn rely on a mixture of fourier analytic and algebraic methods (aside: the mathematical tendency to consider also the machinery of integrability is not 
universally  %generally 
appropriate  %here 
- most physically interesting models are not integrable, while most do benefit from translation symmetry/fourier and spectral/representation theoretic methods). 
Thus it is, for example, that XXZ spin chains, 2d Potts models, six-vertex models and so on have algebras such as Temperley--Lieb algebras \cite{TemperleyLieb71} in their computational machinery. 
More recently motivations such as topological quantum computation have driven interest in topological phases, and statistical mechanical models of systems with boundary doping 
(see e.g. \cite{Fradkin,GuidrySun} and references therein).  
Just as for homogeneous systems, one can consider the algebraic layer in these settings. There are various interesting examples, but here 
in this paper 
we focus on the construction given in \cite{KadarYu} - hence `Kadar--Yu' algebras. 
In this construction there is a tower of algebras $J_l = J_l(\xx)$ for each integer $l \geq -1$;
and then for each value of the parameter $\xx$ taken from the ground field. 
(Here the default ground field will be $\C$.) 
In statistical mechanical terms the index $l$ sets how far the doping will penetrate into the bulk, but controlled topologically rather than metrically. In algebraic terms it controls which subset of a set of Brauer generators are present.
%...

From the representation theory perspective the KY algebras are also intrinsically interesting. We have already mentioned Temperley--Lieb, and the Brauer algebras \cite{Brauer} and partition algebras \cite{Martin91} also arise both as transfer matrix algebras 
(see e.g. \cite{Gainutdinov_2014,Ramgoolam07,Ramgoolam} and references therein)
and objects of direct study in representation theory. 
For example, the non-semisimple representation theories of Temperley--Lieb and Brauer algebras can be seen as controlled by appropriate different cases of Kazhdan--Lusztig theory. The resultant theories are very different in terms of their fundamental invariants, and so the KY algebras provide an opportunity to interpolate between these cases, and hence 
study  %learn about 
geometrical representation theory itself, as we will see. 
%...

In this paper we 
%\bajm{[unfinished... can I help or is this best left to you, Paul?]} \ppm{[I can do it (am doing it), but if you have some pithy phrases then of course drop them in!]}
%...
focus on representation theory. 
A key technical result is Theorem~\ref{thm:main}.
In principle this reduces the computation of 
`bootstrap' 
gram matrices for each tower $J_l$ 
to the computation of one fixed and two initial polynomials per $\l\vdash l+2$, 
all of which are then routine to determine. 
We then illustrate this with a comprehensive set of examples (see \S\ref{ss:examples0}). 
The Theorem and examples raise two key questions: 
firstly how to characterise roots (analogous to the characterisation for classical 
Chebyshev in terms of roots of unity) - see \S\ref{s:roots};
and then how to run the Tower of Recollement programme from this bootstrap to obtain representation theory
- see \S\ref{ss:repthy}. 
For this in particular we need also in principle to determine, in the cases where a `one-cup' standard module {\em has} a submodule (flagged by a vanishing determinant), 
what is the structure of this submodule. 
For the algebras we study here we show, in \S\ref{ss:bootcase0}-\ref{ss:keybootl1}, that 
in almost all cases 
\ppmm{no module can map in} except the no-cup module with the same labelling partition; and then that this 
is a composition factor with multiplicity 1. 

An intriguing question is how to determine not just the bootstrap determinants but all standard determinants. 
Indeed this was one of the initial aims of the paper. 
This is important here 
for various reasons. Firstly 
because these algebras $J_{l,n}$ exhibit a large variety of homological complexities, and so 
provide a good laboratory to understand the passage from gram determinants to Cartan decomposition matrices, 
including homological grading (manifestations of Loewy depth of factors in modules and so on). 
Ab initio the determinants are, collectively, a lot of information, and it is a challenge even to start to present this in an accessible way - this is itself interesting, raising the question of how global limit properties of representation theory unify local linear data in the fibre over the limit.  
In the $l=-1$ (Temperley--Lieb) case there is indeed a striking property that enables one to determine 
not just the bootstrap determinants, but all standard determinants (we review it in \S\ref{ss:gdetdecrolletl-1}). 
This leads us to consider {\em marginal vertex functions} in the general case. 
And indeed one then observes empirically that the presentation of results is hugely simplified. 
And furthermore the results boil down to expressions in terms of the same polynomials as before. 
In \S\ref{ss:gdetdecrollet} we review the $l=-1$ case and summarize the data showing how this generalises. 
\\
(Aside: Since we have the bootstrap determinants, another approach in principle is to reverse-engineer other gram determinants from the representation theory determined via the tower of recollement - potentially linking directly to generalised Kazhdan--Lusztig theory. To achieve the overall aim the ultimate challenge would be to understand the problem approached from {\em both} directions, however we will demote this level of completeness to  a later work.)

\medskip 

Another very intriguing set of questions arises from technical aspects of our computations. 
Key to these computations are properties of Chebyshev series here called `\ramping'. 
These are important for representation theory since they determine properties of the (discrete but rich) subset of parameter values for non-semisimplicity. 
They thus generalise 
(necessarily quite profoundly)
the `root of unity' phenomenon from classical cases. This is a key part of the aimed-for interpolation between Temperley--Lieb-like and Brauer-like representation theories. But the Chebyshev aspect is intrinsically interesting, as we shall see in \S\ref{s:roots}. 
%...

\medskip 


\subsection{Glossary of notations and terms}

See Table~\ref{tab:glossary}. 

\newcommand{\Brp}{{\mathcal B}}   %% hook for Brauer partitions %%

\begin{table}[]
    \centering
    \begin{tabular}{l|l|l}
    notation & short description & where  \\  
     & & defined \\ \hline
    $\Brp(n,m) $ & full set of $(n,m)$ Brauer partitions  &  \ref{de:Brp} \\ 
    $p \mapsto p^*$ & natural map $*:\Brp(n,m) \rightarrow \Brp(m,n)$ (diagram vertical flip) & \ref{de:flip} \\ 
    $J_{l,n}$   &  Kadar--Yu algebra at rank $n$ and height $l$ & \ref{de:Jln} \\  
    $J_l$     &  Kadar--Yu category at height $l$  &   \ref{de:Jln} \\ 
     $\Deltaaa{n}{(p,\l)} $   &  standard module of $J_{l,n}$ for fixed $l$ and index $(p,\l)$ & \ref{de:stdmods} \\
%     \sout{$ \Delt{n}{x}  $} &  \sout{gram matrix of $\Deltaaa{n}{x}$ wrt given basis} & \ref{de:grammat} \\ 
     $ \Gamma(\Delt{n}{p,\l})  $ &  gram matrix of $\Deltaaa{n}{p,\l}$ wrt given basis & \ref{de:gramdet} \\ 
     $\Delttt{(n)}{x}$  & gram determinant of $ \Gamma(\Delt{n}{p,\l})  $ wrt given basis &\ref{de:gramdet} \\ 
     $\Lambda_n$ & set of integer partitions of $n$ &  \ref{de:Lambdan}  \\ 
     $\Lambda_{l,n}$ & index set for standard modules of $J_{l,n}$  & \ref{de:labelset}  \\ 
     $\Lambda_{l,\infty}$ & index set for standard modules in global/large $n$ limit  & \ref{def:indexset}  \\ 
     %$\Sp_\l$  &  &  \\ 
     $\SS_\l$ & Specht module of $\Sym_n$ (for some $n$) with index $\l \in \Lambda_n$  &  \ref{eq:Specht0} \\ 
     $\Sym_n$ & symmetric group of rank $n$\\
     $\uu_{ij}$ & fix $n$. cup (at $\{i,j\}$) element of $\Brp(n,n-2)$  & \ref{de:uu} \\ 
     $\uu_{ij} \uu_{i'j'}^*$ & fix $n$. cup-cap (at $\{i,j\},\{ i',j' \}$) element of $\Brp(n,n)$  & \ref{de:uu} \\ 
    \end{tabular}
    \caption{Glossary of notations %\ppm{ [maybe comment if/when some notation different from \cite{KadarYu,ShBr}?]}
    } %\bajm{[check refs]}}
    \label{tab:glossary}
\end{table}


\newpage 

